% !TeX root = ../main.tex
% 摘要

\ustcsetup{
  keywords  = {多维度代码覆盖率, 嵌入式软件, 异构平台支持, 覆盖率融合},
  keywords* = {Multidimensional code coverage, Embedded software, Heterogeneous platform support, Coverage data fusion},
}

\begin{abstract}
    随着嵌入式软件系统复杂度不断增加，传统覆盖率采集与分析方法难以满足多平台、多领域的测试需求。特别是在资源受限的异构嵌入式平台中，缺乏独立文件系统和测试后自动断电等限制，对覆盖率的实时、可靠采集构成了挑战。此外，传统的全量覆盖率检视方法耗时且容易出错。因此，设计一个支持异构平台覆盖率采集，能够高效整合和可视化呈现多维度覆盖率数据的系统显得尤为重要。

    系统分为三个主要功能模块：覆盖率采集，细粒度覆盖率解析，覆盖率可视化。覆盖率采集模块中，集成了共享内存通信机制，解决了无独立文件系统环境下覆盖率数据实时传输的难题。细粒度覆盖率解析模块中，设计实现了通用的自动化解析流程，兼容主流的交叉编译工具链（LLVM 和 GCC），并通过模块-文件映射和代码差异分析，实现对全量覆盖率数据的细粒度拆分和增量覆盖率提取，突破了传统粗粒度覆盖率检视的限制，提升了测试人员定位盲点的效率。覆盖率可视化模块中，针对多领域、多时间维度的覆盖率数据，设计了融合与可视化方案，分层次、多维度、细粒度地呈现全面、直观的覆盖率数据视图。最终，通过将覆盖率分析脚本部署至 Smart-CI 持续集成平台，实现了覆盖率采集与分析任务的自动化、分布式执行，显著提升了测试效率和资源利用率。

    在实际应用中，该系统已在部门内部得到推广并获得积极反馈，充分验证了其工程可行性与价值，为嵌入式软件多维度覆盖率管理提供了有效的解决方案与实践参考。
\end{abstract}

\begin{abstract*}
    With the increasing complexity of embedded software systems, traditional coverage collection and analysis methods are becoming inadequate for testing needs across multiple platforms and domains. In particular, resource-constrained heterogeneous embedded platforms — characterized by the absence of an independent file system and automatic power-off after test execution — pose significant challenges to real-time and reliable coverage data acquisition. Furthermore, conventional full-coverage inspection methods are time-consuming and error-prone. These challenges highlight the need for a system that supports coverage collection on heterogeneous platforms while efficiently integrating and visualizing multidimensional coverage data.

    The proposed system is organized into three main functional modules: coverage collection, fine-grained coverage analysis, and coverage visualization. In the coverage collection module, a shared-memory communication mechanism is integrated to address the difficulty of real-time coverage data transfer in environments without an independent file system. In the fine-grained coverage analysis module, a general-purpose automated parsing process is implemented, compatible with mainstream cross-compilation toolchains (LLVM and GCC). Through module-to-file mapping and code difference analysis, the system enables fine-grained decomposition of full coverage data and extraction of incremental coverage, thereby overcoming the limitations of traditional coarse-grained inspection and improving testers’ ability to identify blind spots. In the coverage visualization module, a data fusion and visualization scheme is designed for multi-domain and multi-temporal coverage data, presenting comprehensive, hierarchical, and fine-grained views. Finally, by deploying the coverage analysis scripts on the Smart-CI continuous integration platform, the system achieves automated and distributed execution of coverage collection and analysis tasks, significantly enhancing testing efficiency and resource utilization.

    In practical applications, the system has been adopted within the department and received positive feedback, demonstrating its engineering feasibility and value. This work provides an effective solution and practical reference for multidimensional coverage management in embedded software testing.

\end{abstract*}
